% چکیده فارسی؛ حداکثر ۶۰۰ کلمه و مطابق بند ۲-۸ دستورالعمل دانشگاه
\ChapterStart
\begin{center}
  {\XTitr\fontsize{16pt}{20pt}\selectfont \ThesisTitleFA\par}
  \vspace{1.2em}
  {\XTitr\fontsize{16pt}{20pt}\selectfont چکیده}
\end{center}
\vspace{1em}
\noindent\textbf{مقدمه و هدف:}
بیماران پس از انفارکتوس میوکارد (\lr{MI})، با وجود دریافت درمان‌های دارویی و مداخلات بازگشایی عروق، همچنان در معرض عود حوادث قلبی‌عروقی قرار دارند. رعایت اقدامات غیردارویی پیشگیری ثانویه، شامل پرهیز از دخانیات، فعالیت بدنی کافی، الگوی غذایی محافظ قلب و مدیریت مناسب وضعیت روان‌شناختی، برای کاهش این خطر ضروری است. بااین‌حال، پایبندی هم‌زمان به این رفتارها در مراقبت روزمره با چالش‌های متعددی روبه‌رو است. این مطالعه با هدف تعیین میزان پایبندی به اقدامات غیردارویی پیشگیری ثانویه و شناسایی عوامل جمعیت‌شناختی و بالینی مرتبط با آن در بیماران مبتلا به انفارکتوس میوکارد انجام شد.

\vspace{0.6em}
\noindent\textbf{روش کار:}
این پژوهش مقطعی‌-‌تحلیلی در نیمه دوم سال ۱۴۰۴ بر روی ۱۵۰ بیمار ۱۸ سال و بالاتر انجام شد که دست‌کم یک ماه از وقوع \lr{MI} آنان گذشته بود و برای پیگیری به درمانگاه قلب بیمارستان شهید مدنی خرم‌آباد مراجعه کردند. نمونه‌گیری به‌صورت تمام‌شماری و گردآوری داده‌ها با مصاحبه حضوری و تلفنی انجام گرفت. اطلاعات جمعیت‌شناختی و بالینی ثبت شد و مصرف دخانیات، فعالیت بدنی، پایبندی به رژیم غذایی مدیترانه‌ای، وضعیت روان‌شناختی، سواد سلامت، ادراک بیماری و پایبندی دارویی با ابزارهای استاندارد سنجیده شدند. پایبندی کامل به‌معنای تحقق هم‌زمان معیارهای چهار حیطه غیردارویی بود. داده‌ها با آمار توصیفی، آزمون‌های مقایسه‌ای و رگرسیون لجستیک دوحالتی تحلیل شدند و سطح معنی‌داری کمتر از ۰٫۰۵ در نظر گرفته شد.

\vspace{0.6em}
\noindent\textbf{یافته‌ها:}
میانگین سن بیماران $62.1 \pm 10.7$ سال بود و ۷۰ درصد آنان مرد بودند. ۸۲ درصد بیماران در طی 30 روز قبل از ارزیابی، مصرف دخانیات نداشتند؛ ۶۹٫۳ درصد فعالیت بدنی کافی داشتند؛ ۶۰ درصد به رژیم غذایی مدیترانه‌ای پایبند بودند؛ و وضعیت روان‌شناختی ۴۶ درصد طبیعی بود. تنها ۲۶ بیمار، معادل ۱۷٫۳ درصد، به‌طور هم‌زمان به همه اقدامات غیردارویی پایبند بودند. هیچ بیماری سابقه شرکت در برنامه بازتوانی قلبی نداشت. در تحلیل چندمتغیره، مردان و افراد شاغل با احتمال بیشتری مصرف سیگار را گزارش کردند. جنسیت مرد، اشتغال و نوع درمان با پایبندی به فعالیت بدنی ارتباط مستقل داشتند. وضعیت اقتصادی متوسط/خوب، احتمال پایبندی به رژیم غذایی را افزایش داد. تأهل با وضعیت روان‌شناختی بهتر همراه بود و در مقابل، زن بودن و تصور منفی‌تر از بیماری با دیسترس روان‌شناختی ارتباط داشتند. تحصیلات دانشگاهی با پایبندی کامل ارتباط معنی‌دار داشت.

\vspace{0.6em}
\noindent\textbf{نتیجه‌گیری:}
پایبندی کامل به مجموعه اقدامات غیردارویی پیشگیری ثانویه در بیماران مورد مطالعه پایین بود و عوامل مرتبط با هر حیطه یکسان نبودند. این یافته‌ها بر ضرورت ارزیابی جداگانه رفتارهای پیشگیرانه، حمایت هدفمند از بیماران دارای تحصیلات پایین‌تر یا محدودیت اقتصادی، توجه ویژه به فعالیت بدنی زنان، غربالگری وضعیت روان‌شناختی و ایجاد مسیر دسترسی به برنامه بازتوانی قلبی تأکید می‌کنند. با توجه به طراحی مقطعی و تک‌مرکزی مطالعه، نتایج بیانگر ارتباط آماری هستند و نباید به‌عنوان رابطه علت و معلولی تفسیر شوند.

\vspace{1em}
\noindent\textbf{کلمات کلیدی:} \KeywordsFA

\clearpage
